\section{First Letter}

\begin{quotationx}
\textbf{Julius Evola} first met the young \textbf{Mircea Eliade} during a trip to Bucharest to meet Corneliu Codreanu\footnote{\url{https://gornahoor.net/?p=4303}}. Evola refers to that meeting decades later in his spiritual autobiography, \emph{The Cinnabar Path}. By that time, Eliade had become a widely respected historian of religions in the West and kept his early affiliations a secret. Henceforth, Eliade cut off contact with Evola, which they had maintained sporadically until the 1960s.

Nevertheless, Eliade kept up with all of Evola's (and Guenon's) writings as they were published. We translated Eliade's review of Revolt against the Modern World\footnote{\url{https://gornahoor.net/?p=4303}}.

In this letter Evola makes the rather surprising admission that he had been considering moving permanently to India to undergo some spiritual practice, believing he had accomplished all he could in the West. Evola brings up an important point which should be heeded by all. The most authentic centers in India are remote, difficult to find, and closed to foreigners. India today is replete with spiritual frauds and charlatans which makes it extremely difficult for Westerners to encounter a fully traditional teaching; besides, they are prone to be ``gamed'', much like Margaret Mead. Those Hindus who come to the West\footnote{\url{http://www.religioperennis.org/documents/rcoomaraswy/HINDU.PDF}} are often suspect.
\end{quotationx}


28 May 1930

Rome

I received your letter. I remember you perfectly. One of your friends here in Rome had already told me that you had gone to India in November 1928 .

I would very much like to know what you found there in the natural order of things that are of interest to us: that of practice, more than that of doctrine or metaphysics.

I was thinking, and am still thinking (since I am at the point of having finished what I had attempted to do in the West) of going to India to stay there. One of my correspondents convinced me that it would not be worth the trouble, unless I go to Kashmir or Tibet and I have a way to introduce myself into some of the rarest centers that still conserve the Tradition but are excessively suspicious of any foreigners.

Nevertheless, I would be grateful if you could inform me of what you found in addition, with the understanding not from the cultural or metaphysical point of view.

I am sending you:

\begin{itemize}
\item One of the last existing copies of the complete \emph{Ur} collection from 1928.
\item The complete \emph{Krur} collection from 1929.
\item My book on Tantra Man as power 
\end{itemize}
My books that have appeared since then are:

\begin{itemize}
\item \emph{Pagan Imperialism}
\item \emph{Theory of the Absolute Individual}
\item \emph{Phenomenology of the Absolute Individual}
\end{itemize}


The last two constitute the systematic and definitive exposition of my doctrine. Currently, I am editor of ``La Torre'', two issues of which I include. I didn't edit any journals before \emph{Ur}.

Aside from what you received, there is only the Ur collection of 1927 which is out of print. If you want, I can let you know if there is anyone who can sell it and at what price.

If professor Dasgupta\footnote{\url{https://en.wikipedia.org/wiki/Surendranath\_Dasgupta}}, with whom you stayed January to September 1930, is the author of the books on Hindu philosophy, please ask him if Sir Douglas Ainslie\footnote{\url{https://en.wikipedia.org/wiki/Douglas\_Ainslie}}, whom he knows very well and is my friend, remembered to write to him---as he had promised me---so that he has his publisher send me the two volumes that have already appeared, which I could talk about in Italy or Germany.

\flrightit{Posted on 2012-09-03 by Cologero}

\begin{center}* * *\end{center}

\begin{footnotesize}
\begin{sffamily}
\texttt{cosmodromium on 2012-09-03 at 09:32 said:}

Interesting to read Dasgupta was also familiar with Evola.

Thank You.

\hfill

\texttt{Mihai on 2012-09-04 at 05:51 said:}

I agree with Evola.

Besides, I really believe that if one were to practice a Hindu path, it would be impossible to do so without actually living, permanently, in India.

Although the possibility of actually coming into contact with an authentic spiritual center are very dim, while the possibility for frauds is immense, especially given the guilability of modern westerners.

The chances were scant in the 1930s. Imagine how much so they are today.

\hfill

\texttt{Avery on 2012-09-04 at 11:43 said:}

I also agree with Evola. Recently we have seen the chief editor of a certain publishing company move to India to live with the Gaudiya Vaishnavas. If that group can perform a Guenonian initiation they are doing a funny job showing it.

\hfill

\texttt{Cologero on 2012-09-04 at 12:53 said:}

One likely reason, Avery, is that the Gaudiya Vaishnava is open to foreigners and not restricted to Indians of a specific caste.

\hfill

\texttt{fraktal on 2013-04-01 at 18:15 said:}

Try to PODCAST this article :

The difficult encounter in Rome. Mircea Eliade's post-war relation with Julius Evola--new letters and data--

because we couldn' t ... .

\hfill

\texttt{Mark Citadel on 2015-04-13 at 17:31 said:}

I'm sure I'm not the only one who would kill for a transcript of Evola and Codreanu's discussion together.The ascetic meets the hero. Unfortunately, I highly doubt one was taken. Lost to the passage of history now, since both men are dead.

\hfill

\texttt{Cologero on 2015-04-14 at 21:48 said:}

While not exactly a transcript, Evola wrote about his meeting with Codreanu in an article published in La Vita Italiana in 1938. Don't know if there is much interest in it (and I believe variations of it are available online).

You may also be interested in Evola's impressions of \textit{José Antonio Primo de Rivera}\footnote{\url{http://www.gornahoor.net/?page\_id=1202}}. Evola recognizes the spiritual impulse common to both Codreanu and Primo de Rivera, while rejecting it in his own case.

\hfill

\texttt{Mark Citadel on 2015-04-16 at 17:50 said:}

I'll be sure to check that out! Codreanu fascinates me, and I'm just starting to familiarize myself with Primo De Rivera.

Yes, Evola's article on the meeting was called `The Tragedy of Romania's Iron Guard' or something to that effect. In it he speaks very highly of the Capitanul, who seems to have spellbound him, as he did a Hungarian Jewish historian whom I can't remember the name of, but who met Codreanu as a young boy. It seems fair to call the man a Reactionary martyr. Tragic arc of struggle for unflinching principles.


\hfill

\end{sffamily}
\end{footnotesize}

